\documentclass[11pt,a4paper]{article}

\usepackage[T1]{fontenc}
\usepackage[utf8]{inputenc}
\usepackage[polish]{babel}
\usepackage[none]{hyphenat}
\usepackage{lmodern}
\usepackage{microtype}
\usepackage{geometry}
\usepackage{booktabs}
\usepackage{tabularx}
\usepackage{longtable}
\usepackage{array}
\usepackage{graphicx}
\usepackage{hyperref}
\usepackage{xcolor}
\usepackage{enumitem}
\usepackage{titlesec}
\usepackage{fancyhdr}
\usepackage{csquotes}
\usepackage{caption}
\usepackage{float}

\geometry{margin=2.4cm}
\setlength{\emergencystretch}{3em}
\hyphenpenalty=10000
\exhyphenpenalty=10000
\hypersetup{
  colorlinks=true,
  linkcolor=blue!50!black,
  citecolor=blue!50!black,
  urlcolor=blue!60!black
}

\pagestyle{fancy}
\fancyhf{}
\fancyhead[L]{\small Przetwarzanie obrazów w biometrii Projekt}
\fancyhead[R]{\small Detekcja makroskładników ze zdjęcia}
\fancyfoot[C]{\thepage}
\renewcommand{\headrulewidth}{0.4pt}

\title{%
  Detekcja makroskładników jedzenia na podstawie zdjęcia:\\
  przegląd rozwiązań, modeli i stanu badań%
}
\author{%
  Michał Dąbrowski \and Dominik Zabielski\\[0.4em]
  \normalsize Przetwarzanie obrazów w biometrii Projekt%
}
\date{październik 2026}

\begin{document}
\maketitle

\begin{abstract}
Dokument stanowi przegląd aktualnego stanu wiedzy i narzędzi w obszarze
automatycznego szacowania wartości odżywczych (kalorie oraz makroskładniki:
białko, tłuszcz, węglowodany) na podstawie fotografii posiłku.
Omówiono główne sformułowania problemu, kluczowe zbiory danych
(w szczególności Nutrition5k), klasyczne i współczesne architekturę
(CNN, Vision Transformer, multimodalne VLM), wybrane modele dostępne na
Hugging Face oraz rozwiązania komercyjne (aplikacje mobilne).
Na końcu wskazano praktyczne ścieżki realizacji projektu proseminaryjnego.
\end{abstract}

\tableofcontents
\newpage

%% =====================================================================
\section{Wprowadzenie i motywacja}
%% =====================================================================

Ręczne logowanie posiłków w aplikacjach dietetycznych jest czasochłonne
i podatne na błędy. Alternatywą jest \emph{image-based dietary assessment}:
system, który na podstawie zdjęcia talerza szacuje masę porcji oraz
zawartość kalorii i makroskładników.

Problem jest trudny z kilku powodów:
\begin{itemize}[leftmargin=*]
  \item z jednego 2D zdjęcia nie da się jednoznacznie odtworzyć objętości
        (głębia, okluzja, perspektywa);
  \item potrawy są złożone (wiele składników, sosy, warstwy);
  \item ta sama klasa jedzenia ma bardzo różną gęstość kaloryczną
        ponadto niektóre składniki mogą być niewidoczne na zdjęciu, na przykład
        oliwa w sałatce;
  \item bazy żywieniowe różnią się między krajami (USDA, krajowe FCDB).
\end{itemize}

Mimo to dziedzina dynamicznie się rozwija: od klasycznej klasyfikacji
Food-101, przez regresję makro na Nutrition5k (Google, CVPR~2021),
po modele multimodalne (VLM) z dostępem do baz USDA / FNDDS.

%% =====================================================================
\section{Sformułowanie problemu}
%% =====================================================================

W praktyce stosowane są trzy główne podejścia.

\subsection{Podejście A: klasyfikacja i baza odżywcza}
\begin{enumerate}[leftmargin=*]
  \item Rozpoznaj klasę potrawy / składniki (klasyfikacja lub detekcja).
  \item Oszacuj porcję (heurystyka, referencyjny obiekt, głębokość).
  \item Zmapuj wynik do bazy (USDA FoodData Central, Open Food Facts itp.).
\end{enumerate}
Zaleta: interpretowalność. Wada: błąd kumuluje się na każdym etapie;
słaba generalizacja na dania mieszane.

\subsection{Podejście B: bezpośrednia regresja makroskładników}
Sieć (CNN / ViT) dostaje obraz i przewiduje wektor
\[
  \hat{y} = (\mathrm{kcal},\; m,\; \mathrm{protein},\; \mathrm{fat},\; \mathrm{carbs}).
\]
(odpowiednio: kalorie, masa, białko, tłuszcz, węglowodany).
To dominujące podejście naukowe na Nutrition5k.
Zaleta: prosty trening od początku do końca. Wada: słabsza interpretowalność;
trudniej przenieść model na inną kuchnię / bazę.

\subsection{Podejście C: multimodalny VLM i wyszukiwanie w bazie}
Model wizyjno-językowy opisuje składniki i gramy, a wartości odżywcze
pobierane są z autorytatywnej bazy (np.\ DietAI24 + FNDDS).
Zaleta: szerszy zakres składników odżywczych, lepsza elastyczność.
Wada: koszt inferencji, ryzyko halucynacji bez RAG.

\textbf{Najbardziej sensowne podejście na ten moment:} rozpocząć od B,
czyli regresji na zbiorze Nutrition5k. Następnie warto dodać A, czyli
klasyfikację Food-101 jako punkt odniesienia. Podejście C należy potraktować
jako wariant zaawansowany, możliwy do wykonania dopiero po uzyskaniu
działającej wersji podstawowej.

%% =====================================================================
\section{Kluczowe zbiory danych}
%% =====================================================================

\begin{table}[H]
\centering
\caption{Wybrane zbiory istotne dla tematu.}
\label{tab:datasets}
\begin{tabular}{@{}llp{6.2cm}@{}}
\toprule
\textbf{Zbiór} & \textbf{Skala} & \textbf{Co daje} \\
\midrule
Nutrition5k & $\sim$5006 dań & RGB(+D), wideo, masa, kcal, makro, składniki (USDA) \\
Food-101 & 101 klas $\times$ 1000 & klasyfikacja potraw (bez dokładnych makro) \\
FoodSeg103 & segmentacja & maski składników (pomoc przy porcjach) \\
NutritionVerse & syntetyczne + rzeczywiste & wielozadaniowa regresja makro \\
Open Food Facts & produkty opakowane & etykiety / kody kreskowe (nie talerz) \\
\bottomrule
\end{tabular}
\end{table}

\subsection{Nutrition5k (Google Research, CVPR 2021)}
Najważniejszy zbiór odniesienia do bezpośredniej estymacji makro:
\begin{itemize}[leftmargin=*]
  \item $\sim$5000 prawdziwych dań z kantyn Google;
  \item $>250$ składników, średnio $\sim$5{,}7 składnika na talerz;
  \item etykiety: całkowita masa, kalorie, tłuszcz, węglowodany, białko
        oraz rozbicie per-składnik;
  \item dla $\sim$3500 dań dostępne obrazy RGB-D z góry (Intel RealSense);
  \item oficjalny podział na zbiór uczący i testowy ($\sim$90/10).
\end{itemize}
Źródło: \url{https://github.com/google-research-datasets/Nutrition5k}.
Na Hugging Face istnieją przetworzone warianty, między innymi:
\texttt{ai5labsOfficial/nutrition5k},\\
\texttt{arunapb/nutrition5k-foodseg103}.

\textbf{Uwaga praktyczna:} pełny zbiór jest duży ($\sim$180\,GB).
Do eksperymentów studenckich wystarczy podzbiór klatek RGB rejestrowanych z góry
+ metadane CSV.

%% =====================================================================
\section{Architektury i podejścia naukowe (stan 2024--2026)}
%% =====================================================================

\subsection{Model bazowy na Nutrition5k}
Oryginalna praca Thames i~in.\ (CVPR~2021) pokazała, że regresja z obrazu
może przewyższać szacunki profesjonalnych dietetyków na tym zbiorze.
Dodanie głębokości (RGB-D) poprawia estymację masy/porcji.

\subsection{CNN i Vision Transformers}
Benchmarki 2025 (m.in.\ analiza na Nutrition5k z różnymi FCDB)
wskazują, że:
\begin{itemize}[leftmargin=*]
  \item \textbf{ResNet-101} dobrze radzi sobie z masą i energią;
  \item \textbf{ViT-B/16} często wygrywa na białku i węglowodanach;
  \item InceptionV3 (historyczny model bazowy) jest już wyraźnie słabszy.
\end{itemize}
Typowe błędy (MAPE, rząd wielkości): masa $\sim$20--23\%,
energia $\sim$26--29\%, makro $\sim$30--50\%, zależnie od przygotowania danych.
Korekta mas składników w etykietach mocno poprawia metryki (nawet 6--42\%).

\textbf{NutritionVerse-Direct} (2024): ViT + wspólne warstwy FC + 5 głowic
regresyjnych; istotna poprawa względem Inception-ResNet.

\textbf{Swin Transformer + EfficientNet + FPN} (Frontiers in Nutrition, 2024):
Top-1 $\sim$79{,}5\%, MAPE kalorii $\sim$14{,}7\% na Nutrition5k.

\textbf{Cross-ViT wieloskalowy} (2026): dla samego obrazu RGB raportowano
średni błąd bezwzględny około 28\,kcal oraz około 6\,g białka, 5\,g tłuszczu
i 10\,g węglowodanów.

\subsection{Segmentacja + open-vocabulary}
Podejścia łączące segmentację jedzenia (oddzielenie tła) z Transformerem
i/lub głębokością osiągają średnie PMAE rzędu $\sim$17\% na Nutrition5k.
Analizy mechanizmów uwagi pokazują, że tło bywa ważne dla kalorii/masy/węglowodanów,
a mniej dla białka i tłuszczu.

\subsection{Estymacja głębi na podstawie RGB}
Gdy nie ma czujnika głębi, stosuje się modele estymacji głębi
(np.\ Depth Anything) i dokłada kanał głębokości. Przykładem jest
NutriIngredientNet / DPF-Nutrition.

\subsection{Multimodalne LLM i VLM z wyszukiwaniem w bazie}
\textbf{DietAI24} (Communications Medicine, 2025): MLLM + Retrieval-Augmented
Generation z bazą FNDDS; szacuje dziesiątki składników odżywczych,
nie tylko makroskładniki. Silnie redukuje średni błąd bezwzględny
względem klasycznych modeli komputerowego widzenia
na mieszanych daniach rzeczywistych.

Trend: \emph{nie polegać na ``pamięci'' VLM co do kcal}, tylko
rozpoznawać składniki/gramy i czytać bazę.

%% =====================================================================
\section{Modele na Hugging Face: przegląd praktyczny}
%% =====================================================================

Poniżej modele bezpośrednio przydatne do projektu (stan na październik 2026).
Liczby pobrań zmieniają się szybko, dlatego ważniejsza jest jakość karty modelu
i zgodność z Nutrition5k / makro.

\subsection{Estymacja kalorii / makro (priorytet tematyczny)}

\scriptsize
\begin{longtable}{@{}p{5.8cm}p{2.4cm}p{6.0cm}@{}}
\toprule
\textbf{Model (HF)} & \textbf{Typ} & \textbf{Komentarz} \\
\midrule
\endhead
\texttt{opencal/opencal-base} & VLM (ONNX) &
  Model dostrojony na bazie LFM2.5-VL-450M; JSON składników z gramami
  oraz kcal i makroskładnikami (USDA); działa w przeglądarce
  (transformers.js / WebGPU). \\
\texttt{jc-builds/CalorieCLIP} & CLIP + regresja &
  MAE $\sim$51{,}4\,kcal; lokalnie, $<$50\,ms; tylko kalorie (nie pełne makro);
  lepszy od GPT-4o/Claude w ich teście. \\
\texttt{aman24012/nutrivision-om-igsm} & CNN/hybryda &
  5 wyjść: cal, mass, fat, carb, prot; mean PMAE $\sim$13{,}5\% na Nutrition5k. \\
\texttt{arunapb/nutriingredientnet-v5} & RGB + głębia &
  Per-składnik + makro; MAE kcal $\sim$79{,}7; wymaga geometrii zbliżonej
  do ujęć z góry ze zbioru Nutrition5k. \\
\texttt{Doses-AI/\allowbreak boba-0.8b-food-GGUF} & mały VLM &
  Qwen3.5-0.8B LoRA; JSON dla każdego składnika; MAE kcal $\sim$112;
  działa na urządzeniu (GGUF). \\
\texttt{ardaocak/llava-1.5-7b-food-\allowbreak calorie-estimator} & LLaVA &
  Estymacja kalorii ze zdjęcia; cięższy, wymaga GPU. \\
\bottomrule
\end{longtable}
\normalsize

\subsection{Klasyfikacja potraw (model bazowy, podejście A)}

\begin{itemize}[leftmargin=*]
  \item \texttt{onnx-community/swin-finetuned-food101-ONNX}: popularny Swin na Food-101;
  \item \texttt{prithivMLmods/Food-101-93M}: ViT i klasyfikator Food-101;
  \item \texttt{nateraw/vit-base-food101}: klasyczny ViT-B na Food-101;
  \item \texttt{Kaludi/food-category-classification-v2.0}: szersze kategorie żywności.
\end{itemize}
Te modele \textbf{nie} przewidują makro bezpośrednio. Trzeba dociągnąć
bazę odżywczą i model porcji.

\subsection{Modele tekstowe / pomocnicze}
Na HF jest też sporo LoRA do LLM (SmolLM2, Phi-4, Llama-3, Gemma)
uczących się mapować opis jedzenia $\rightarrow$ wartości odżywcze
(np.\ rodzina \texttt{punkt2/foodcalculator-nutrition-*}).
Przydatne, gdy najpierw VLM opisze talerz tekstem.

\subsection{Jak szybko zacząć na HF}
\begin{enumerate}[leftmargin=*]
  \item Demonstrator internetowy: CalorieCLIP lub NutriVision (jeśli dostępny).
  \item Lokalnie: \texttt{transformers} + \texttt{opencal-base} (ONNX)
        albo własna głowica regresji na ViT/ResNet pretrained.
  \item Trening: podzbiór Nutrition5k i uczenie wielozadaniowe
        z funkcją błędu MSE lub Huber dla 5 wyjść.
\end{enumerate}

%% =====================================================================
\section{Rozwiązania komercyjne i produktowe}
%% =====================================================================

Rynek aplikacji ``zrób zdjęcie posiłku'' jest dojrzały produktowo,
ale dokładność bywa mocno marketingowa.

\begin{itemize}[leftmargin=*]
  \item \textbf{MyFitnessPal}: ogromna baza produktów; Meal Scan (analiza zdjęcia)
        jako dodatek do klasycznego logowania / kodów kreskowych.
  \item \textbf{Foodvisor}: rozpoznawanie ze zdjęcia oraz makro- i mikroskładników;
        mocna pozycja w Europie.
  \item \textbf{SnapCalorie / Cal AI / podobne}: aplikacje skupione na analizie zdjęć;
        niezależne benchmarki 2026 pokazują dużą rozpiętość MAPE między aplikacjami;
        dania mieszane są najtrudniejsze.
  \item \textbf{Cronometer / MacroFactor}: często wygrywają dokładnością
        przy ręcznym wprowadzaniu danych i kodach kreskowych
        (nie przy samej analizie zdjęcia).
\end{itemize}

Wniosek: aplikacje konsumenckie łączą CV z ogromnymi bazami i UX
(korekta użytkownika). Sam model ze zdjęcia rzadko jest wystarczająco
dokładny bez weryfikacji.

%% =====================================================================
\section{Metryki ewaluacji}
%% =====================================================================

Standardowo raportuje się osobno dla kalorii, masy i każdego makro:
\begin{itemize}[leftmargin=*]
  \item MAE / RMSE (jednostki bezwzględne: kcal, g),
  \item MAPE / PMAE (błąd względny w \%),
  \item czasem korelację Pearsona / Bland--Altman.
\end{itemize}
Przy uczeniu wielozadaniowym warto podawać też średnią po 5 celach
oraz unikać porównań między publikacjami bez tego samego podziału danych
i tego samego sposobu przygotowania etykiet.

%% =====================================================================
\section{Wyzwania i ograniczenia}
%% =====================================================================

\begin{enumerate}[leftmargin=*]
  \item \textbf{Porcja z 2D}: największe źródło błędu; głębia lub obiekt referencyjny
        (karta kredytowa, sztućce) pomaga.
  \item \textbf{Zmiana dziedziny}: model z kantyny Google $\neq$ domowa kuchnia polska.
  \item \textbf{Dania złożone}: sałatki, dania jajeczne, śniadania zachodnie
        są systematycznie trudne.
  \item \textbf{Gęste kalorycznie, małe objętościowo}: orzechy, oleje:
        modele często niedoszacowują.
  \item \textbf{Baza odżywcza}: zmiana USDA $\leftrightarrow$ krajowe tabele
        zmienia wartości referencyjne.
  \item \textbf{Ograniczenia VLM}: lokalne VLM bywają gorsze od małej
        regresji CLIP (jak w CalorieCLIP).
\end{enumerate}

%% =====================================================================
\section{Propozycja realizacji projektu (dla grupy)}
%% =====================================================================

Cel projektu: system przyjmujący zdjęcie posiłku i zwracający
szacunki \textbf{kalorii + białka + tłuszczu + węglowodanów}
(opcjonalnie masa).

\subsection{MVP (realistyczny zakres projektu)}
\begin{enumerate}[leftmargin=*]
  \item Model bazowy klasyfikacji Food-101 (gotowy model HF) i prosta tabela kcal/makro.
  \item Regresja wielozadaniowa na podzbiorze Nutrition5k:
        backbone ViT-B/16 lub ResNet-50 + 5 głowic.
  \item Porównanie z gotowym modelem HF (CalorieCLIP dla kcal;
        NutriVision / OpenCal na makro jeśli uda się odpalić).
  \item Prosty interfejs: notebook lub mała aplikacja (Streamlit / Gradio).
  \item Raport błędów i analiza przypadków (np.\ dania jedno- i wieloskładnikowe).
\end{enumerate}

\subsection{Rozszerzenia (jeśli zostanie czas)}
\begin{itemize}[leftmargin=*]
  \item estymacja głębi z RGB i połączenie kanałów;
  \item segmentacja jedzenia vs tło;
  \item VLM i wyszukiwanie w bazie USDA / Open Food Facts;
  \item test na własnych zdjęciach kuchni polskiej i ręczne oznaczenie danych.
\end{itemize}

\subsection{Technologie (propozycja)}
Python, PyTorch, Hugging Face \texttt{transformers} / \texttt{timm},
opcjonalnie ONNX Runtime, Gradio; metryki: MAE, MAPE.

%% =====================================================================
\section{Podsumowanie}
%% =====================================================================

Estymacja makroskładników ze zdjęcia to aktywna dziedzina na styku
komputerowego widzenia i analizy żywienia. Najważniejszą podstawą naukową
jest \textbf{Nutrition5k} oraz regresja wielozadaniowa (CNN/ViT).
Na Hugging Face dostępne są zarówno lekkie regresory (CalorieCLIP),
jak i VLM produkujące JSON ze składnikami i makro (OpenCal, Boba).
Komercyjnie dominują aplikacje łączące analizę zdjęć z bazami i korektą użytkownika.

W naszym projekcie chcemy skupić się przede wszystkim na zbiorze
\textbf{Nutrition5k} i modelu ViT albo ResNet uczonym jednocześnie do
przewidywania kilku wartości. Taki model będzie podstawą naszych
eksperymentów. Dla porównania sprawdzimy też jeden lub dwa gotowe modele
z Hugging Face i przygotujemy prosty demonstrator działania.

%% =====================================================================
\section*{Wybrane źródła i linki}
\addcontentsline{toc}{section}{Wybrane źródła i linki}
%% =====================================================================

\begin{enumerate}[leftmargin=*,itemsep=0.35em]
  \item Thames i~in., \emph{Nutrition5k: Towards Automatic Nutritional Understanding of Generic Food}, CVPR 2021.
        \url{https://github.com/google-research-datasets/Nutrition5k}
  \item NutritionVerse-Direct (2024): \url{https://arxiv.org/abs/2405.07814}
  \item 2D Prediction\ldots Beyond Nutrition5k, \emph{Nutrients} 2025:
        \url{https://doi.org/10.3390/nu17132196}
  \item Swin Transformer nutritional composition, \emph{Frontiers in Nutrition} 2024:
        \url{https://doi.org/10.3389/fnut.2024.1454466}
  \item DietAI24, \emph{Communications Medicine} 2025:
        \url{https://doi.org/10.1038/s43856-025-01159-0}
  \item HF: \url{https://huggingface.co/opencal/opencal-base}
  \item HF: \url{https://huggingface.co/jc-builds/CalorieCLIP}
  \item HF: \url{https://huggingface.co/aman24012/nutrivision-om-igsm}
  \item HF: \url{https://huggingface.co/arunapb/nutriingredientnet-v5}
  \item HF: \url{https://huggingface.co/Doses-AI/boba-0.8b-food-GGUF}
  \item USDA FoodData Central: \url{https://fdc.nal.usda.gov/}
\end{enumerate}


\end{document}
